\providecommand{\Res}[2]{\qty[round-mode=figures,round-precision=3]{\fpeval{#1}}{#2}}
\providecommand{\Val}[1]{\num[round-mode=figures,round-precision=4]{\fpeval{#1}}}
\providecommand{\ResSci}[2]{\qty[round-mode=figures,round-precision=3,exponent-mode=scientific]{\fpeval{#1}}{#2}}
\subsubsection*{Question a) Angle de diffusion}
\begin{enumerate}[label=\arabic*)]
    \item \textbf{Variables connues et recherchées :} $\lambda = \qty{\v{l0}}{pm}$, $\lambda' = \qty[round-mode=places,round-precision=3]{\fpeval{\v{l0} + 1e12*6.626e-34/(9.109e-31*3e8)*(1 - cosd(\v{th}))}}{pm}$, $h/(m_e c) = \qty{2.425}{pm}$ ; variable recherchée : $\theta$
    \item \textbf{Équations pertinentes :}
    \[ \lambda' - \lambda = \frac{h}{m_e c}(1 - \cos\theta) \]
    \item \textbf{Isolation de la variable recherchée :}
    \[ \theta = \arccos\left(1 - \frac{(\lambda' - \lambda)\,m_e c}{h}\right) \]
    \item \textbf{Résolution numérique :}
    \[ \Delta\lambda = \Res{1e12*6.626e-34/(9.109e-31*3e8)*(1 - cosd(\v{th}))}{pm} \qquad \theta = \arccos\left(1 - \frac{\Val{1e12*6.626e-34/(9.109e-31*3e8)*(1 - cosd(\v{th}))}}{\num{2.425}}\right) = \Res{\v{th}}{\degree} \]
    \item \textbf{Réponse :} le détecteur est placé à \textbf{\Res{\v{th}}{\degree}}.
\end{enumerate}

\subsubsection*{Question b) Énergie de l'électron}
\[ K_e = hc\left(\frac{1}{\lambda} - \frac{1}{\lambda'}\right) = \ResSci{6.626e-34*3e8*(1/\vsi{l0} - 1/(\vsi{l0} + 6.626e-34/(9.109e-31*3e8)*(1 - cosd(\v{th}))))}{J} = \mathbf{\Res{6.626e-34*3e8*(1/\vsi{l0} - 1/(\vsi{l0} + 6.626e-34/(9.109e-31*3e8)*(1 - cosd(\v{th}))))/1.602e-19}{eV}} \]

\subsubsection*{Question c) Photons non décalés}
Ils ont été diffusés par des électrons fortement liés aux atomes : c'est alors tout l'atome, beaucoup plus massif qu'un électron, qui recule. Le décalage $h/(Mc)$ est négligeable.
